\section{The row walk and its occupations}
\label{sec:flows}

This section decides how often to apply each split and each tripling of
Section~\ref{sec:maps}. Apart from the copies of \(S^+_{a,a}\) that are
fed in and the copies of \(\MM_2\) that come out, every tensor factor that
one relation produces must be consumed by another relation, and
conversely; in addition, the logarithms of the auxiliary dimensions,
weighted by how often each split is used, must have a small sum, because
the auxiliary dimensions multiply into the source side of the final
scheme. We make these choices with a random walk in each row, the
\emph{row walk}, whose expected numbers of visits, the
\emph{occupations}, balance every factor exactly.
Proposition~\ref{prop:occupations} proves this and bounds both the total
occupation and the weighted sum of logarithms by
\(\Gamma_a=105+180(a-2)\).

\subsection{States, splits and triplings}
\label{subsec:row-walk}

Fix the target index \(a\ge2\) and a row \(r\) with \(2\le r\le a\), and
put \(\sigma_r=(-1)^{a-r}\). Row \(r\) consists of the tensors
\(S^{\sigma_r}_{r,b}\) with columns \(r-1\le b\le4r-1\). We index them by
the \emph{state} \(x=b-r+1\in\{0,\ldots,3r\}\), so that state \(x\) is
column \(b=r-1+x\). With \(f_r(x)=\frac{3(r-1)}2+x\), the row-affine
form~\eqref{eq:row-affine} of the least profile reads, in the state
variable,
\begin{equation}\label{eq:affine-row}
 F(r,r-1+x)=s_rf_r(x)\qquad(0\le x\le3r).
\end{equation}

The states \(1\le x\le3r-1\) are \emph{interior}. For an interior state
\(x\) we define its \emph{width}, its auxiliary dimension, its
\emph{children} and its transition probabilities by
\begin{align}
 h_{r,x}&=\min(x,3r-x),&d_{r,x}&=h_{r,x}+1,\notag\\
 \mathcal Y_{r,x}&=\{x-h_{r,x},x-h_{r,x}+2,\ldots,x+h_{r,x}\},
 &q_{r,x,y}&=\frac{f_r(y)}{d_{r,x}f_r(x)}
 \quad(y\in\mathcal Y_{r,x}).\label{eq:multiway-probabilities}
\end{align}
The \(d_{r,x}\) children are placed symmetrically about \(x\), so their
average is \(x\). Since \(f_r\) is affine, it follows that
\(\sum_{y\in\mathcal Y_{r,x}}f_r(y)=d_{r,x}f_r(x)\); hence the
\(q_{r,x,y}\) are positive rationals with sum one. By
\eqref{eq:affine-row}, a child is chosen with probability proportional to
the profile value \(F(r,r-1+y)\) at its column; in other words,
\((h+1)F(r,b)=\sum_{j=0}^hF(r,b+h-2j)\) for \(b=r-1+x\) and
\(h=h_{r,x}\). This makes \(\log F\) a potential that telescopes along
the walk (Proposition~\ref{prop:entropy-loss}).

Since \(h_{r,x}=\min(x,3r-x)\) is the largest \(h\) with \(0\le x-h\) and
\(x+h\le3r\), every child satisfies \(0\le y\le3r\), so its column
\(r-1+y\) lies in \([r-1,4r-1]\), and no wider split keeps all its
children in the row. This maximal width makes the walk cross a row
quickly, which keeps the cost of a row bounded
(Proposition~\ref{prop:occupations}). Since \(r\ge2\) and
\(h_{r,x}\le x\), the column of \(x\) satisfies
\(b=r-1+x\ge h_{r,x}+1\), which is the hypothesis of
Lemma~\ref{lem:multiway}.

The walk keeps track of tensor factors. State \(x\) of row \(r\) stands
for the factor \(S^{\sigma_r}_{r,r-1+x}\). A step from an interior state
\(x\) stands for one \emph{split} at \((r,x)\): the restriction of
Lemma~\ref{lem:multiway} with \(b=r-1+x\), \(h=h_{r,x}\) and auxiliary
\(B_X(d_{r,x})\), used in the three orientations of
\(\mathcal A_{\sigma_r}\) as in Corollary~\ref{cor:oriented-relations},
so that the three auxiliaries combine into \(\MM_{d_{r,x}}\)
by~\eqref{eq:coset-dot-products}. Its block with filtration index
\(j=(x+h_{r,x}-y)/2\) is \(S^{\sigma_r}_{r,r-1+y}\), the factor of the
child \(y\). In a tensor power, each copy of the factor at \(x\) becomes
one factor at some child, and the type of the restriction fixes how many
copies go to each child; the walk moves from \(x\) to \(y\) with the
frequency \(q_{r,x,y}\) that this type approximates
(Section~\ref{sec:rounding}).

The two end states have special roles. State \(3r\) is column \(4r-1\),
and the step from it to state \(1\), column \(r\), stands for one
\emph{tripling}: the restriction of Lemma~\ref{lem:tripling} in the three
orientations of \(\mathcal A_{\sigma_r}\), which turns copies of
\(S^{\sigma_r}_{r,4r-1}\) into sectors that each contain the same number
of copies of \(S^{\sigma_r}_{r,r}\), by~\eqref{eq:cyclic-tripling-balance}.
State \(0\) is column \(r-1\). By~\eqref{eq:cyclic-identities} and
\(\sigma_{r-1}=-\sigma_r\), its factor
\(S^{\sigma_r}_{r,r-1}\cong S^{\sigma_{r-1}}_{r-1,r}\) is the factor at
state \(2\) of row \(r-1\) (Section~\ref{subsec:orientations}). A walk
that reaches state \(0\) therefore leaves row \(r\) and continues in row
\(r-1\) at state \(2\). Hence the walk of a lower row starts at state
\(2\), whereas the walk of the top row starts at state \(1\), the diagonal
factor \(S^{\sigma_a}_{a,a}=S^+_{a,a}\) whose powers are the input of the
construction; accordingly, the \emph{injection state} of row \(r\) is
\(\xin_r=1\) if \(r=a\) and \(\xin_r=2\) if \(r<a\). In row \(2\), state
\(0\) is
\(S^{\sigma_2}_{2,1}\cong S^{-\sigma_2}_{1,2}\cong\MM_2\), the output of
the construction.

\begin{definition}[Row walk and occupations]\label{def:row-walk}
The \emph{row walk} of row \(r\) is the Markov chain
\(X_0,X_1,X_2,\ldots\) on the states \(\{0,\ldots,3r\}\) with
\(X_0=\xin_r\) and the following transitions. From an interior state
\(x\), it moves to the child \(y\in\mathcal Y_{r,x}\) with probability
\(q_{r,x,y}\); from state \(3r\), it moves to state \(1\); and state \(0\)
is absorbing. The \emph{occupation} \(\lambda_{r,x}\) of an interior
state \(x\) is the expected number of times \(\nu\ge0\) with
\(X_\nu=x\), that is, the expected number of splits at \((r,x)\). The
expected number of times \(\nu\) with \(X_\nu=3r\) is denoted by
\(\mu_r\); it is the expected number of triplings in row \(r\).
\end{definition}

A priori these expectations lie in \([0,\infty]\);
Proposition~\ref{prop:occupations} shows that they are finite and
rational. Figure~\ref{fig:row-walk} and Example~\ref{ex:row-two} show the
row walk of row \(2\).

\begin{figure}[ht]
\centering
\begin{tikzpicture}[
  x=2.25cm, y=1cm, >={Stealth[length=2mm]},
  st/.style={circle,draw,thick,minimum size=7.5mm,inner sep=0pt},
  un/.style={st,draw=black!45,text=black!55},
  pr/.style={font=\footnotesize,inner sep=1pt,fill=white},
  gpr/.style={pr,text=black!60},
  tr/.style={->,thick},
  gr/.style={->,draw=black!40}
  ]
  \node[st,double,double distance=1.2pt] (s0) at (0,0) {\(0\)};
  \node[st] (s1) at (1,0) {\(1\)};
  \node[st] (s2) at (2,0) {\(2\)};
  \node[un] (s3) at (3,0) {\(3\)};
  \node[st] (s4) at (4,0) {\(4\)};
  \node[un] (s5) at (5,0) {\(5\)};
  \node[st] (s6) at (6,0) {\(6\)};
  % moves to the right (above the states)
  \draw[tr] (s1) to[bend left=30] node[pr]{\(7/10\)} (s2);
  \draw[tr] (s2) to[bend left=38] node[pr]{\(11/21\)} (s4);
  \draw[tr] (s4) to[bend left=38] node[pr]{\(5/11\)} (s6);
  \draw[gr] (s3) to[bend left=30] node[gpr]{\(11/36\)} (s4);
  \draw[gr] (s3) to[bend left=45] node[gpr]{\(5/12\)} (s6);
  \draw[gr] (s5) to[bend left=30] node[gpr]{\(15/26\)} (s6);
  % moves to the same state
  \draw[tr] (s2) to[out=115,in=65,looseness=6]
    node[pr,above,fill=none]{\(1/3\)} (s2);
  \draw[tr] (s4) to[out=115,in=65,looseness=6]
    node[pr,above,fill=none]{\(1/3\)} (s4);
  % moves to the left (below the states)
  \draw[tr] (s1) to[bend left=30] node[pr]{\(3/10\)} (s0);
  \draw[tr] (s2) to[bend left=38] node[pr]{\(1/7\)} (s0);
  \draw[tr] (s4) to[bend left=38] node[pr]{\(7/33\)} (s2);
  \draw[gr] (s3) to[bend left=45] node[gpr]{\(1/12\)} (s0);
  \draw[gr] (s3) to[bend left=30] node[gpr]{\(7/36\)} (s2);
  \draw[gr] (s5) to[bend left=30] node[gpr]{\(11/26\)} (s4);
  % tripling and injection
  \draw[->,thick,dashed] (s6) to[out=95,in=85,looseness=0.55]
    node[pr,above,fill=none]{tripling} (s1);
  \draw[->,very thick] (0.55,1.0)
    node[pr,above left,fill=none]{injection if \(a=2\)} -- (s1);
  \draw[->,very thick] (2,-2.0)
    node[pr,below,fill=none]{injection if \(a>2\)} -- (s2);
  % columns
  \node[font=\footnotesize,anchor=east] at (-0.35,0) {state \(x\)};
  \foreach \x/\b in {0/1,1/2,2/3,3/4,4/5,5/6,6/7}
    \node[font=\footnotesize] at (\x,-2.9) {\(\b\)};
  \node[font=\footnotesize,anchor=east] at (-0.35,-2.9) {column \(b\)};
\end{tikzpicture}
\caption{The row walk of row \(r=2\) (Example~\ref{ex:row-two}). Arcs
above the states move right, arcs below move left, and loops stay; labels
are the probabilities \(q_{2,x,y}\). Dashed: the tripling \(6\to1\).
Double circle: the absorbing state \(0\), the factor
\(S^{\sigma_2}_{2,1}\cong\MM_2\). Gray: states not reachable from the
injection states.}
\label{fig:row-walk}
\end{figure}

\begin{example}[Row \(r=2\)]\label{ex:row-two}
Here \(s_2=4/3\) and \(f_2(x)=x+3/2\), so \(F(2,b)=\frac43(b+\frac12)\)
for \(1\le b\le7\), and the states \(0,\ldots,6\) are the columns
\(1,\ldots,7\). The widths of the interior states \(1,\ldots,5\) are
\(1,2,3,2,1\). For instance, the children of state \(2\) are \(0,2,4\);
their values of \(f_2\) are \(3/2,7/2,11/2\), whose sum is
\(3f_2(2)=21/2\), so \(q_{2,2,0}=1/7\), \(q_{2,2,2}=1/3\) and
\(q_{2,2,4}=11/21\). Figure~\ref{fig:row-walk} shows all transition
probabilities. Every child in this row is even: a state \(x\le3\) has the
children \(0,2,\ldots,2x\), and a state \(x\ge3\) has the children
\(2x-6,\ldots,6\). Consequently a walk that starts at state \(1\) or
\(2\) never visits states \(3\) and \(5\), and
\(\lambda_{2,3}=\lambda_{2,5}=0\). For a state other than \(0\), the
expected number of visits equals the number of starts there plus the
expected number of arrivals. For the states \(1,2,4,6\) this gives
\begin{align*}
 \lambda_{2,1}&=\delta_{1,\xin_2}+\mu_2,&
 \lambda_{2,2}&=\delta_{2,\xin_2}+\tfrac7{10}\lambda_{2,1}
   +\tfrac13\lambda_{2,2}+\tfrac7{33}\lambda_{2,4},\\
 \lambda_{2,4}&=\tfrac{11}{21}\lambda_{2,2}+\tfrac13\lambda_{2,4},&
 \mu_2&=\tfrac5{11}\lambda_{2,4},
\end{align*}
where \(\delta\) is the Kronecker delta. These are the coordinates
\(1,2,4,6\) of the row identity~\eqref{eq:row-flow} below.

If \(a=2\), then \(\xin_2=1\), and the solution is \(\lambda_{2,1}=2\),
\(\lambda_{2,2}=14/5\), \(\lambda_{2,4}=11/5\) and \(\mu_2=1\): the walk
makes \(7\) splits on average, and
\(\sum_x\lambda_{2,x}\log d_{2,x}=2\log2+5\log3\approx6.88\). If \(a>2\),
then \(\xin_2=2\), and the solution is \(\lambda_{2,1}=10/7\),
\(\lambda_{2,2}=4\), \(\lambda_{2,4}=22/7\) and \(\mu_2=10/7\): the walk
makes \(60/7\approx8.57\) splits on average, and
\(\sum_x\lambda_{2,x}\log d_{2,x}=\frac{10}7\log2+\frac{50}7\log3
\approx8.84\). In both cases \(\mu_2\) agrees
with~\eqref{eq:reset-weights}.
\end{example}

\subsection{Occupations}
\label{subsec:occupations}

Let \(\mathbf e_0,\ldots,\mathbf e_{3r}\) be the standard
basis of \(\Q^{3r+1}\), indexed by the states of row \(r\).

\begin{proposition}[Occupations]\label{prop:occupations}
Let \(a\ge2\) and \(2\le r\le a\). The row walk of row \(r\), started at
any state, is absorbed at state \(0\) in finite expected time. Its
occupations
\(\lambda_{r,x}\), for \(1\le x\le3r-1\), and \(\mu_r\) are nonnegative
rationals. They are the unique solution of a nonsingular rational linear
system, and they satisfy the row identity
\begin{equation}\label{eq:row-flow}
 \sum_{x=1}^{3r-1}\lambda_{r,x}
 \Bigl(\mathbf e_x-\sum_{y\in\mathcal Y_{r,x}}q_{r,x,y}\mathbf e_y\Bigr)
 +\mu_r(\mathbf e_{3r}-\mathbf e_1)=\mathbf e_{\xin_r}-\mathbf e_0.
\end{equation}
The expected numbers of triplings are
\begin{equation}\label{eq:reset-weights}
 \mu_a=\frac1{a-1},\qquad
 \mu_r=\frac{2(3r-1)}{(r-1)(3r+1)}\quad(r<a);
\end{equation}
in particular, \(\mu_r\le2\). Each of the two row sums
\(\sum_x\lambda_{r,x}\) and \(\sum_x\lambda_{r,x}\log d_{r,x}\) is less
than \(105\) if \(r=a\) and less than \(180\) if \(r<a\). Consequently,
each \(\lambda_{r,x}\) and each \(\lambda_{r,x}\log d_{r,x}\) is less
than \(180\), and, with \(\Gamma_a=105+180(a-2)\),
\begin{equation}\label{eq:occupation-bounds}
 \sum_{r=2}^a\sum_{x=1}^{3r-1}\lambda_{r,x}\le\Gamma_a,
 \qquad
 \sum_{r=2}^a\sum_{x=1}^{3r-1}\lambda_{r,x}\log d_{r,x}\le\Gamma_a.
\end{equation}
\end{proposition}

On the left side of~\eqref{eq:row-flow}, the coefficient of
\(\mathbf e_y\) is the expected number of steps out of \(y\) minus the
expected number of steps into \(y\); by the right side, it is \(1\) at
\(\xin_r\), \(-1\) at state \(0\) and \(0\) elsewhere, so one unit enters
the row at \(\xin_r\) and leaves it at state \(0\). In terms of tensor
factors, the splits and triplings of row \(r\), used with these frequencies, consume
one more copy of \(S^{\sigma_r}_{r,r-1+\xin_r}\) than they produce,
produce one more copy of \(S^{\sigma_r}_{r,r-1}\) than they consume, and
balance every other factor of the row.

\begin{proof}
The idea is to compare each excursion of the row walk between visits to
the end states \(0\) and \(3r\) with the \emph{uniform walk}, which picks
every child with equal probability. The uniform walk is a martingale, and
the function \(g\) below decreases in expectation by at least
\(\gamma\sqrt{h_{r,x}}\) at each step from \(x\); this bounds the
expected number of splits in an excursion and the expected sum of
\(\log d_{r,x}\) over them. Between visits to the end states, the row
walk is the Doob transform of the uniform walk by the affine function
\(f_r\) \citep[Section~17.6.1]{LevinPeres17}; this changes these
expectations by at most the factor \(f_r(3r)/f_r(1)=3\). Finally,
triplings chain the excursions together, and their number is
geometrically distributed.

\emph{Notation.} Fix the row and drop the index \(r\) from \(f_r\),
\(h_{r,x}\), \(d_{r,x}\), \(\mathcal Y_{r,x}\), \(q_{r,x,y}\),
\(\lambda_{r,x}\), \(\mu_r\) and \(\xin_r\). Let
\(\theta_{\mathrm{ex}}=\min\{\nu\ge0:X_\nu\in\{0,3r\}\}\) be the exit
time from the interior (\(\theta_{\mathrm{ex}}=\infty\) if there is no
such \(\nu\)). We write \(\mathbb P^q_x\) and \(\mathbb E^q_x\) for the
row walk started at \(x\), and \(\mathbb P^{\mathrm u}_x\) and
\(\mathbb E^{\mathrm u}_x\) for the uniform walk started at \(x\), which
moves from an interior state \(x\) to each child \(y\in\mathcal Y_x\)
with probability \(1/d_x\) and stops at \(0\) and at \(3r\). An
\emph{excursion} of the row walk runs from an \emph{entry} state (the
starting state, or state \(1\) right after a tripling) up to the next
visit to \(0\) or \(3r\). A \emph{charge} is a function \(\varphi\ge0\)
on the interior states, and the charge of a path is
\(\sum_{\nu<\theta_{\mathrm{ex}}}\varphi(X_\nu)\). We use \(\varphi=1\),
which counts splits, and \(\varphi(x)=\log d_x\). Both satisfy
\(\varphi(x)\le\sqrt{h_x}\) at interior states, the first because
\(h_x\ge1\), and the second because \(\log(1+h)\le\sqrt h\) for
\(h\ge1\): the difference \(\sqrt h-\log(1+h)\) is positive at \(h=1\),
and its derivative \((\sqrt h-1)^2/\bigl(2\sqrt h\,(1+h)\bigr)\) is
nonnegative.

\emph{Step 1: a square-root drift.} Put
\(g(x)=\sqrt x+\sqrt{3r-x}-\sqrt{3r}\) and \(\gamma=1-2\sqrt2/3>0\).
Since the square root is subadditive, \(g\ge0\) on \([0,3r]\). We claim
that
\begin{equation}\label{eq:square-root-drift}
 g(x)-\mathbb E^{\mathrm u}_x\,g(X_1)\ge\gamma\sqrt{h_x}
 \qquad(1\le x\le3r-1).
\end{equation}
First let \(x\le3r/2\). Then \(h_x=x\), and the children are
\(0,2,\ldots,2x\), each with probability \(1/(x+1)\). On each interval
\([\ell-1,\ell]\), the concave square root lies above its chord, so
\[
 \sum_{\ell=0}^x\sqrt\ell-\frac{\sqrt x}2
 =\sum_{\ell=1}^x\frac{\sqrt{\ell-1}+\sqrt\ell}2
 \le\int_0^x\sqrt\xi\,d\xi=\frac23x^{3/2}.
\]
Dividing by \(x+1\) and using \(\frac23x+\frac12\le\frac23(x+1)\), we
obtain
\[
 \mathbb E^{\mathrm u}_x\sqrt{X_1}
 =\frac{\sqrt2}{x+1}\sum_{\ell=0}^x\sqrt\ell
 \le\sqrt2\,\sqrt x\;\frac{\frac23x+\frac12}{x+1}
 \le\frac{2\sqrt2}3\sqrt x.
\]
The children average to \(x\), so Jensen's inequality for the concave
function \(\xi\mapsto\sqrt{3r-\xi}\) gives
\(\mathbb E^{\mathrm u}_x\sqrt{3r-X_1}\le\sqrt{3r-x}\). Together,
\[
 g(x)-\mathbb E^{\mathrm u}_x\,g(X_1)
 =\bigl(\sqrt x-\mathbb E^{\mathrm u}_x\sqrt{X_1}\bigr)
  +\bigl(\sqrt{3r-x}-\mathbb E^{\mathrm u}_x\sqrt{3r-X_1}\bigr)
 \ge\Bigl(1-\frac{2\sqrt2}3\Bigr)\sqrt x=\gamma\sqrt{h_x}.
\]
For \(x\ge3r/2\), the reflection \(\xi\mapsto3r-\xi\) maps
\(\mathcal Y_x\) onto \(\mathcal Y_{3r-x}\), preserves the width, and
leaves \(g\) unchanged. Hence the first case, applied to the state
\(3r-x\le3r/2\), proves~\eqref{eq:square-root-drift} at \(x\).

\emph{Step 2: excursions of the uniform walk.} We show that
\begin{equation}\label{eq:uniform-excursion-bound}
 \mathbb E^{\mathrm u}_x\sum_{\nu<\theta_{\mathrm{ex}}}\sqrt{h_{X_\nu}}
 \le\frac{g(x)}\gamma\le\frac{\sqrt x}\gamma
 \qquad(1\le x\le3r-1).
\end{equation}
Fix an integer \(\bar\nu\ge1\). On the event
\(\{\nu<\theta_{\mathrm{ex}}\}\) the state \(X_\nu\) is interior, so the
Markov property and~\eqref{eq:square-root-drift} give
\(\mathbb E^{\mathrm u}_x\bigl[g(X_\nu)-g(X_{\nu+1})\bigm|X_0,\ldots,X_\nu\bigr]
\ge\gamma\sqrt{h_{X_\nu}}\) there. The event
\(\{\nu<\theta_{\mathrm{ex}}\wedge\bar\nu\}\) is determined by
\(X_0,\ldots,X_\nu\). Multiplying by its indicator, taking expectations
and summing over \(\nu<\bar\nu\), we obtain
\[
 \gamma\,\mathbb E^{\mathrm u}_x\sum_{\nu<\theta_{\mathrm{ex}}\wedge\bar\nu}\sqrt{h_{X_\nu}}
 \le\mathbb E^{\mathrm u}_x\sum_{\nu<\theta_{\mathrm{ex}}\wedge\bar\nu}
   \bigl(g(X_\nu)-g(X_{\nu+1})\bigr)
 =g(x)-\mathbb E^{\mathrm u}_x\,g(X_{\theta_{\mathrm{ex}}\wedge\bar\nu})\le g(x),
\]
where the last step uses \(g\ge0\). Letting \(\bar\nu\to\infty\),
monotone convergence gives the first inequality
in~\eqref{eq:uniform-excursion-bound}, and \(\sqrt{3r-x}\le\sqrt{3r}\)
gives the second. This argument does not assume that
\(\theta_{\mathrm{ex}}\) is finite; since \(1\le\sqrt{h_x}\) and
\(\varphi(x)\le\sqrt{h_x}\) at interior states, it gives
\(\mathbb E^{\mathrm u}_x\theta_{\mathrm{ex}}\le\sqrt x/\gamma<\infty\)
and \(\mathbb E^{\mathrm u}_x\sum_{\nu<\theta_{\mathrm{ex}}}\varphi(X_\nu)\le\sqrt x/\gamma\)
for both charges.

\emph{Step 3: excursions of the row walk.} The function \(f\) is positive,
since \(f(0)=3(r-1)/2>0\), and harmonic for the uniform walk at every
interior state, because the children average to \(x\) and \(f\) is
affine. By~\eqref{eq:multiway-probabilities},
\(q_{x,y}=d_x^{-1}f(y)/f(x)\) is the Doob transform of the uniform
probabilities by \(f\). The tripling step \(3r\to1\) is not of this form,
since \(f(1)\ne f(3r)\), so we transform only excursions: for every path
\(x_0,x_1,\ldots,x_\nu\) whose states \(x_0,\ldots,x_{\nu-1}\) are
interior, whether or not \(x_\nu\) is interior, the ratio of the path
probabilities telescopes:
\[
 \mathbb P^q_{x_0}(X_1=x_1,\ldots,X_\nu=x_\nu)
 =\prod_{\ell=1}^\nu\frac{f(x_\ell)}{d_{x_{\ell-1}}f(x_{\ell-1})}
 =\frac{f(x_\nu)}{f(x_0)}\,
  \mathbb P^{\mathrm u}_{x_0}(X_1=x_1,\ldots,X_\nu=x_\nu).
\]
Since \(f\) is increasing, every entry state \(x_0\ge1\) and every state
\(x_\nu\le3r\) satisfy
\(f(x_\nu)/f(x_0)\le f(3r)/f(1)=3\), because \(f(1)=(3r-1)/2\) and
\(f(3r)=3(3r-1)/2\). For each \(\nu\), the event
\(\{\nu<\theta_{\mathrm{ex}}\}\) consists of the paths whose states \(X_0,\ldots,X_\nu\) are all interior. Summing the
path identity over these paths with \(x_0=x\) and applying the last bound
gives
\(\mathbb E^q_x\bigl[\varphi(X_\nu)\mathbf 1_{\{\nu<\theta_{\mathrm{ex}}\}}\bigr]
\le3\,\mathbb E^{\mathrm u}_x\bigl[\varphi(X_\nu)\mathbf 1_{\{\nu<\theta_{\mathrm{ex}}\}}\bigr]\).
We sum over \(\nu\) and use Step~2. The expected charge
\(\psi_x=\mathbb E^q_x\sum_{\nu<\theta_{\mathrm{ex}}}\varphi(X_\nu)\) of
an excursion of the row walk from the entry state \(x\) therefore
satisfies
\begin{equation}\label{eq:harmonic-excursion-bound}
 \psi_x\le\frac{3\sqrt x}\gamma\qquad(1\le x\le3r-1)
\end{equation}
for both charges. With \(\varphi=1\), this gives
\(\mathbb E^q_x\theta_{\mathrm{ex}}\le3\sqrt x/\gamma\), so every
excursion of the row walk ends in finite expected time.

\emph{Step 4: exit probabilities.} Under \(\mathbb P^{\mathrm u}_x\), the
stopped process \(X_{\theta_{\mathrm{ex}}\wedge\nu}\) is a martingale with
values in \([0,3r]\), since the children of an interior state average to
that state. By Step~2, \(\theta_{\mathrm{ex}}<\infty\) almost surely, so
bounded convergence gives
\(\mathbb E^{\mathrm u}_xX_{\theta_{\mathrm{ex}}}=x\). As
\(X_{\theta_{\mathrm{ex}}}\in\{0,3r\}\), the uniform walk leaves the
interior at \(3r\) with probability \(x/(3r)\). Summing the path identity
of Step~3 over the paths that first leave the interior at \(3r\), we find
that the row walk started at \(x\) leaves the interior at \(3r\) with
probability
\begin{equation}\label{eq:exit-probability}
 \rho_x=\mathbb P^q_x(X_{\theta_{\mathrm{ex}}}=3r)=\frac{xf(3r)}{3rf(x)},\qquad
 \rho_1=\frac1r,\qquad
 \rho_2=\frac{2(3r-1)}{r(3r+1)}<\frac2r,
\end{equation}
where we used \(f(2)=(3r+1)/2\) and the values of \(f(1)\) and \(f(3r)\)
from Step~3.

\emph{Step 5: chaining the excursions.} Start the row walk at an interior
state \(x\). Each excursion ends either at \(0\), where the walk is
absorbed, or at \(3r\), after which a tripling moves the walk to state
\(1\) and a new excursion starts there. By the strong Markov property at
the exit times, each new excursion is independent of the past and has the
law of an excursion under \(\mathbb P^q_1\); in particular, it ends at
\(3r\) with probability \(\rho_1\). Hence the walk makes at least \(\ell\)
triplings with probability \(\rho_x\rho_1^{\ell-1}\) for \(\ell\ge1\), so
the expected number of triplings is \(\rho_x/(1-\rho_1)\). The excursion
that follows the \(\ell\)-th tripling takes place exactly when at least
\(\ell\) triplings occur, and its expected charge is then \(\psi_1\).
Summing over \(\ell\), we find that the expected total charge of the row
walk started at \(x\) is \(\psi_x+\rho_x\psi_1/(1-\rho_1)\).

Now start at \(x=\xin\), so that the expected number of triplings is
\(\mu\) and the expected total charge is
\(\sum_{y=1}^{3r-1}\lambda_y\varphi(y)\). Since \(1-\rho_1=(r-1)/r\),
\eqref{eq:exit-probability} gives \(\mu=\rho_1/(1-\rho_1)=1/(r-1)\) if
\(r=a\), and \(\mu=\rho_2/(1-\rho_1)=2(3r-1)/\bigl((r-1)(3r+1)\bigr)\)
if \(r<a\). This proves~\eqref{eq:reset-weights}, and both values are at
most \(2/(r-1)\le2\). By~\eqref{eq:harmonic-excursion-bound},
\(\psi_1\le3/\gamma\) and \(\psi_2\le3\sqrt2/\gamma\). Using \(r\ge2\),
through \(r/(r-1)\le2\) and \(2/(r-1)\le2\), and
\(1/\gamma=3/(3-2\sqrt2)=9+6\sqrt2\), we find that the expected total
charge is
\[
 \begin{cases}
  \psi_1+\dfrac{\rho_1\psi_1}{1-\rho_1}=\dfrac r{r-1}\,\psi_1
  \le\dfrac6\gamma=54+36\sqrt2<105,&r=a,\\[10pt]
  \psi_2+\dfrac{\rho_2\psi_1}{1-\rho_1}\le\psi_2+\dfrac2{r-1}\,\psi_1
  \le\dfrac{3\sqrt2+6}\gamma=90+63\sqrt2<180,&r<a.
 \end{cases}
\]
Applied to the charges \(\varphi=1\) and \(\varphi=\log d\), these
estimates prove the bounds on the two row sums.

\emph{Step 6: absorption and the row identity.} Every visit to a state
other than \(0\) is followed by exactly one step. Hence the absorption
time \(\theta_{\mathrm{abs}}=\min\{\nu:X_\nu=0\}\) of the row walk
started at an interior state \(x\) is the number of visits to interior
states plus the number of triplings, and by Step~5 with \(\varphi=1\) its
expectation is \(\psi_x+(1+\psi_1)\rho_x/(1-\rho_1)<\infty\). From state
\(3r\) the walk first moves to state \(1\), so absorption has finite
expected time from every state. Now start the walk at \(\xin\). The
telescoping sum
\(\sum_{\nu<\theta_{\mathrm{abs}}}(\mathbf e_{X_\nu}-\mathbf e_{X_{\nu+1}})
=\mathbf e_{\xin}-\mathbf e_0\) has terms of \(\ell^1\)-norm at most
\(2\), so the norms of its terms add up to at most
\(2\theta_{\mathrm{abs}}\), which has finite expectation. Hence we may
take expectations term by term:
\(\mathbf e_{\xin}-\mathbf e_0=\sum_{\nu\ge0}\mathbb E^q_{\xin}
\bigl[\mathbf 1_{\{\nu<\theta_{\mathrm{abs}}\}}(\mathbf e_{X_\nu}-\mathbf e_{X_{\nu+1}})\bigr]\).
Given \(X_0,\ldots,X_\nu\) with \(X_\nu=x\ne0\), the conditional
expectation of \(\mathbf e_{X_{\nu+1}}\) is
\(\sum_{y\in\mathcal Y_x}q_{x,y}\mathbf e_y\) if \(x\) is interior and
\(\mathbf e_1\) if \(x=3r\). The \(\nu\)-th term is therefore
\[
 \sum_{x=1}^{3r-1}\mathbb P^q_{\xin}(X_\nu=x)
 \Bigl(\mathbf e_x-\sum_{y\in\mathcal Y_x}q_{x,y}\mathbf e_y\Bigr)
 +\mathbb P^q_{\xin}(X_\nu=3r)(\mathbf e_{3r}-\mathbf e_1).
\]
Summing over \(\nu\) and using
\(\sum_\nu\mathbb P^q_{\xin}(X_\nu=x)=\lambda_x\) and
\(\sum_\nu\mathbb P^q_{\xin}(X_\nu=3r)=\mu\), we
obtain~\eqref{eq:row-flow}.

\emph{Step 7: the linear system.} Let \(\Pi\) be the matrix indexed by
the states \(1,\ldots,3r\) with \(\Pi_{x,y}=q_{x,y}\) for interior \(x\)
and \(y\in\mathcal Y_x\setminus\{0\}\), with \(\Pi_{3r,1}=1\), and with
all other entries \(0\). Since state \(0\) is absorbing,
\((\Pi^\nu)_{x,y}=\mathbb P^q_x(X_\nu=y)\) for \(x,y\ge1\), so the row
sums of \(\sum_{\nu\ge0}\Pi^\nu\) are the expected absorption times,
which are finite by Step~6. The series therefore converges, so
\(\Pi^\nu\to0\), and letting \(\bar\nu\to\infty\) in
\(\sum_{\nu=0}^{\bar\nu}\Pi^\nu(\mathrm{Id}-\Pi)=\mathrm{Id}-\Pi^{\bar\nu+1}\)
shows that \(\mathrm{Id}-\Pi\) is nonsingular, with inverse the
fundamental matrix \(\sum_{\nu\ge0}\Pi^\nu\)
\citep[Theorem~11.4]{GrinsteadSnell97}.
The coordinates \(1,\ldots,3r\) of~\eqref{eq:row-flow} state that
\(\boldsymbol\lambda=(\lambda_1,\ldots,\lambda_{3r-1},\mu)\) satisfies
\((\mathrm{Id}-\Pi^{\mathsf T})\boldsymbol\lambda=\mathbf e'_{\xin}\),
where \(\mathbf e'_{\xin}\) is the standard basis vector of \(\Q^{3r}\)
at \(\xin\). This nonsingular system has rational coefficients, so
\(\boldsymbol\lambda\) is its unique solution and is rational, and
Gaussian elimination in exact arithmetic computes it. The occupations are
nonnegative as expected counts.

\emph{Step 8: the totals.} The final row contributes less than \(105\) to
each sum in~\eqref{eq:occupation-bounds}, and each of the \(a-2\) lower
rows contributes less than \(180\). Both sums are therefore at most
\(105+180(a-2)=\Gamma_a\). All terms of the row sums are nonnegative,
since \(d_{r,x}\ge2\). Hence each \(\lambda_{r,x}\) and each
\(\lambda_{r,x}\log d_{r,x}\) is at most its row sum, which is less than
\(180\).
\end{proof}

\begin{remark}[The constants are conservative]\label{rem:occupation-constants}
The constants \(105\) and \(180\) are far above the actual row sums. A
floating-point evaluation of the occupations for \(2\le r\le120\) gives
\(\sum_x\lambda_{r,x}\le7\) and \(\sum_x\lambda_{r,x}\log d_{r,x}<6.88\)
in the final row, and \(\sum_x\lambda_{r,x}\le60/7<8.58\) and
\(\sum_x\lambda_{r,x}\log d_{r,x}<9.98\) in a lower row. The first three
maxima are attained at \(r=2\) (Example~\ref{ex:row-two}), and the last
at \(r=18\).
\end{remark}
